% !TEX root = ../DoAn.tex
\documentclass[../DoAn.tex]{subfiles}
\begin{document}

\chapter{Nền tảng lý thuyết và Công nghệ sử dụng}
\label{chapter:3}

Chương này trình bày các nền tảng lý thuyết và công nghệ được lựa chọn để triển khai hệ thống \textbf{StayJourney}---nền tảng đặt phòng khách sạn trực tuyến đã được phân tích ở Chương~\ref{chapter:2}. Mỗi công nghệ được gắn với yêu cầu chức năng hoặc phi chức năng cụ thể; đồng thời nêu các phương án thay thế và lý do lựa chọn.

\section{Tổng quan ánh xạ yêu cầu--công nghệ}
\label{section:3.1}

Bảng~\ref{tab:yeu-cau-cong-nghe} tóm tắt mối liên hệ giữa nhóm yêu cầu ở Chương~\ref{chapter:2} và công nghệ tương ứng trong đồ án.

\begin{table}[htbp]
\centering
\caption{Ánh xạ nhóm yêu cầu (Chương~\ref{chapter:2}) và công nghệ sử dụng}
\label{tab:yeu-cau-cong-nghe}
\small
\begin{tabular}{|p{4.2cm}|p{5.8cm}|p{4.5cm}|}
\hline
\textbf{Nhóm yêu cầu} & \textbf{Vấn đề cần giải quyết} & \textbf{Công nghệ lựa chọn} \\
\hline
Đặt phòng, giá, khuyến mãi (UC SJ001) & Giao diện tương tác, gọi API, tính giá theo đêm & React, Vite, Axios, Express, MongoDB \\
\hline
Thanh toán QR/VNPay (UC SJ002--SJ004) & Tích hợp cổng thanh toán, lưu minh chứng & VNPay SDK, Multer, Cloudinary \\
\hline
Lưu trú, quản lý đơn (UC SJ005) & Phân quyền theo vai trò, cập nhật trạng thái & RBAC, JWT, Mongoose \\
\hline
Hủy và hoàn tiền (UC SJ006--SJ008) & Kiểm tra chính sách, lưu snapshot & express-validator, service layer \\
\hline
Thông báo thời gian thực & Đẩy sự kiện đến Guest/Owner/Staff & Socket.IO \\
\hline
Bảo mật tài khoản quản trị & Xác thực đa lớp, chống lộ token & JWT (HttpOnly cookie), bcrypt, 2FA email \\
\hline
Giao diện đa thiết bị & Responsive trên desktop/mobile & React, SCSS, Material UI \\
\hline
Vận hành nền & Nhắc check-in, tắt sale hết hạn & node-cron, Nodemailer \\
\hline
\end{tabular}
\end{table}

\section{Kiến trúc hệ thống}
\label{section:3.2}

\subsection{Mô hình client--server và kiến trúc ba tầng}

StayJourney được xây dựng theo mô hình \textbf{client--server}: trình duyệt chạy ứng dụng đơn trang (SPA), máy chủ cung cấp REST API và kết nối cơ sở dữ liệu~\cite{fielding2000rest}. Kiến trúc logic gồm ba tầng: \textit{Presentation} (React), \textit{Business} (Express controllers/services), \textit{Data} (MongoDB qua Mongoose).

\textbf{Yêu cầu liên quan:} toàn bộ nghiệp vụ đa vai trò (Guest, Owner, Staff, Admin) và ba quy trình nghiệp vụ chính (đặt phòng--thanh toán, lưu trú, hủy--hoàn tiền) ở Mục~\ref{subsection:2.2.6}.

\textbf{Phương án thay thế:} kiến trúc monolith server-side rendering (SSR) với template engine; kiến trúc microservices tách từng domain.

\textbf{Lựa chọn:} SPA + REST API vì tách rõ giao diện và nghiệp vụ, dễ triển khai song song hai tầng \texttt{client/} và \texttt{server/}, phù hợp quy mô đồ án và đội phát triển nhỏ. Microservices được loại bỏ do tăng độ phức tạp vận hành mà chưa cần thiết.

\subsection{Giao tiếp REST API}

Các tài nguyên (\texttt{/api/guest}, \texttt{/api/owner}, \texttt{/api/staff}, \texttt{/api/admin}, \texttt{/api/auth}, \texttt{/api/payment}) được thiết kế theo phong cách REST: phương thức HTTP biểu diễn thao tác, JSON làm định dạng trao đổi~\cite{fielding2000rest}.

\textbf{Yêu cầu liên quan:} phân tách API theo vai trò, health check (\texttt{/health}), callback VNPay.

\textbf{Phương án thay thế:} GraphQL; gRPC.

\textbf{Lựa chọn:} REST vì mô hình tài nguyên trực quan, hệ sinh thái Express phong phú, dễ kiểm thử bằng Postman; GraphQL/gRPC phù hợp hơn khi client có nhu cầu truy vấn phức tạp hoặc hiệu năng RPC cao---không phải điểm nghẽn của StayJourney.

\section{Công nghệ phía máy khách (Frontend)}
\label{section:3.3}

\subsection{React và Vite}

\textbf{React} là thư viện JavaScript khai báo giao diện theo component, hỗ trợ tái sử dụng UI và quản lý trạng thái cục bộ~\cite{react2024}. \textbf{Vite} là công cụ build/dev server tận dụng ES modules, rút ngắn thời gian khởi động so với Webpack truyền thống~\cite{vite2024}.

\textbf{Yêu cầu liên quan:} giao diện Guest (tìm KS, đặt phòng), Owner/Staff (sơ đồ phòng, dashboard), Admin (quản trị); yêu cầu responsive.

\textbf{Phương án thay thế:} Angular; Vue.js; Next.js (SSR).

\textbf{Lựa chọn:} React 18 + Vite vì cộng đồng lớn, cú pháp JSX quen thuộc, Vite đủ nhẹ cho đồ án. Next.js SSR không cần thiết vì SEO trang nội bộ quản trị không phải ưu tiên; Angular có đường cong học tập cao hơn.

\subsection{React Router và Redux Toolkit}

\textbf{React Router~6} định tuyến theo URL, hỗ trợ route bảo vệ theo vai trò (\texttt{protected.jsx}). \textbf{Redux Toolkit} quản lý trạng thái toàn cục (phiên đăng nhập, thông báo)~\cite{redux2024}.

\textbf{Yêu cầu liên quan:} bốn không gian \texttt{/}, \texttt{/owner}, \texttt{/staff}, \texttt{/admin}; chuyển hướng sau đăng nhập theo role.

\textbf{Phương án thay thế:} Zustand, Context API thuần; React Router thay bằng TanStack Router.

\textbf{Lựa chọn:} Redux Toolkit cho luồng auth/notifications có nhiều component đồng bộ; React Router là chuẩn de facto cho SPA React.

\subsection{Material UI, SCSS và Axios}

\textbf{Material UI (MUI)} cung cấp component UI theo Material Design, rút ngắn thời gian dựng form, dialog, bảng~\cite{mui2024}. \textbf{SCSS} bổ sung style theo module feature. \textbf{Axios} gọi HTTP kèm cookie (\texttt{withCredentials})~\cite{axios2024}.

\textbf{Yêu cầu liên quan:} giao diện thống nhất, responsive; gọi API có xác thực cookie.

\textbf{Phương án thay thế:} Tailwind CSS + Headless UI; Fetch API thuần.

\textbf{Lựa chọn:} MUI + SCSS kết hợp component sẵn có và tùy biến layout; Axios tiện cấu hình interceptor và cookie hơn Fetch.

\section{Công nghệ phía máy chủ (Backend)}
\label{section:3.4}

\subsection{Node.js và Express}

\textbf{Node.js} chạy JavaScript phía server trên V8, mô hình I/O bất đồng bộ phù hợp API nhiều kết nối~\cite{nodejs2024}. \textbf{Express~4} là framework tối giản định tuyến middleware~\cite{express2024}.

\textbf{Yêu cầu liên quan:} xử lý đồng thời nhiều request đặt phòng, upload, callback thanh toán.

\textbf{Phương án thay thế:} NestJS; Fastify; Spring Boot (Java); Django (Python).

\textbf{Lựa chọn:} Node.js + Express vì cùng ngôn ngữ với frontend, triển khai nhanh, middleware phong phú (\texttt{cors}, \texttt{cookie-parser}). NestJS có cấu trúc hơn nhưng nặng cho phạm vi đồ án.

\subsection{MongoDB và Mongoose}

\textbf{MongoDB} là CSDL NoSQL dạng document BSON, linh hoạt schema~\cite{mongodb2024}. \textbf{Mongoose} cung cấp schema, validation và quan hệ tham chiếu (\texttt{ref})~\cite{mongoose2024}.

\textbf{Yêu cầu liên quan:} lưu đơn đặt phòng (snapshot giá, khuyến mãi), khách sạn, phòng, đánh giá, thông báo; mô hình dữ liệu thay đổi theo chính sách hoàn tiền, sale.

\textbf{Phương án thay thế:} PostgreSQL + Prisma/Sequelize; MySQL.

\textbf{Lựa chọn:} MongoDB vì document JSON khớp payload API, schema linh hoạt cho trường snapshot (\texttt{promotionApplied}, \texttt{nightBreakdown}); quan hệ không quá phức tạp để bắt buộc RDBMS. PostgreSQL phù hợp hơn nếu cần giao dịch ACID phức tạp---StayJourney xử lý nhất quán chủ yếu ở tầng service.

\subsection{express-validator}

Lớp \textbf{validation} tách trong \texttt{server/validations/}, middleware \texttt{validate} trả HTTP~400 kèm danh sách lỗi~\cite{expressvalidator2024}.

\textbf{Yêu cầu liên quan:} kiểm tra đầu vào đặt phòng, tạo KS/phòng, thanh toán, sale---đáp ứng yêu cầu tin cậy dữ liệu.

\textbf{Phương án thay thế:} Joi; Zod; validation thủ công trong controller.

\textbf{Lựa chọn:} express-validator tích hợp sẵn với Express, khai báo rule theo route, tách biệt khỏi logic nghiệp vụ.

\section{Nền tảng bảo mật và phân quyền}
\label{section:3.5}

\subsection{JSON Web Token (JWT) và HttpOnly cookie}

\textbf{JWT} là chuỗi token ký số chứa claims, dùng xác thực stateless giữa các request~\cite{jones2015jwt}. StayJourney dùng cặp \textit{access token} (ngắn hạn) và \textit{refresh token} (dài hạn) lưu trong \textbf{HttpOnly cookie}, giảm rủi ro XSS so với lưu \texttt{localStorage}~\cite{owasp2023session}.

\textbf{Yêu cầu liên quan:} đăng nhập, duy trì phiên, refresh token; bảo vệ API theo role.

\textbf{Phương án thay thế:} session server-side (Redis); OAuth2/OpenID với IdP bên thứ ba.

\textbf{Lựa chọn:} JWT + cookie vì triển khai gọn trên một server Node, không cần Redis; OAuth2 dư thừa khi chưa tích hợp đăng nhập Google/Facebook.

\subsection{bcrypt và xác thực hai lớp (2FA)}

\textbf{bcrypt} băm mật khẩu với salt, chống tấn công brute-force~\cite{provos1999bcrypt}. \textbf{2FA} bắt buộc với Admin/Owner: OTP gửi email, mã dự phòng, danh sách thiết bị tin cậy.

\textbf{Yêu cầu liên quan:} bảo mật tài khoản quản trị và chủ KS; yêu cầu phi chức năng về an toàn thông tin.

\textbf{Phương án thay thế:} Argon2; TOTP (Google Authenticator); SMS OTP.

\textbf{Lựa chọn:} bcrypt đủ cho đồ án, thư viện \texttt{bcryptjs} phổ biến trên Node; 2FA qua email vì không yêu cầu cài app authenticator, phù hợp người dùng khách sạn nhỏ.

\subsection{Phân quyền theo vai trò (RBAC)}

\textbf{RBAC} gán quyền theo \textit{role} thay vì từng user~\cite{sandhu1996rbac}. StayJourney có bốn role; middleware \texttt{requireRole}, \texttt{attachStaffHotel} giới hạn Staff một khách sạn.

\textbf{Yêu cầu liên quan:} tách chức năng Guest/Owner/Staff/Admin ở Chương~\ref{chapter:2}.

\textbf{Phương án thay thế:} ABAC (thuộc tính); ACL từng tài nguyên.

\textbf{Lựa chọn:} RBAC vì mô hình vai trò cố định, dễ mô tả use case; ABAC phức tạp hơn mức cần thiết.

\section{Công nghệ thời gian thực và tác vụ nền}
\label{section:3.6}

\subsection{Socket.IO}

\textbf{Socket.IO} mở rộng WebSocket, hỗ trợ fallback polling, room theo user~\cite{socketio2024}. Server xác thực socket bằng JWT; client React dùng \texttt{socket.io-client} cho chuông thông báo.

\textbf{Yêu cầu liên quan:} thông báo đơn mới, xác nhận thanh toán, phản hồi đánh giá---không cần reload trang.

\textbf{Phương án thay thế:} Server-Sent Events (SSE); polling định kỳ; Firebase Cloud Messaging.

\textbf{Lựa chọn:} Socket.IO vì hai chiều, tích hợp sẵn với HTTP server Express; SSE chỉ một chiều; polling tốn băng thông.

\subsection{node-cron và Nodemailer}

\textbf{node-cron} lên lịch tác vụ (nhắc check-in, vô hiệu sale hết hạn). \textbf{Nodemailer} gửi SMTP: OTP 2FA, hóa đơn email, phản hồi liên hệ~\cite{nodemailer2024}.

\textbf{Yêu cầu liên quan:} email sau thanh toán; nhắc lịch; giao tiếp hỗ trợ.

\textbf{Phương án thay thế:} Bull/BullMQ + Redis; SendGrid/AWS SES.

\textbf{Lựa chọn:} cron trong process Node đủ cho tần suất thấp; Nodemailer + Gmail/SMTP miễn phí cho môi trường demo.

\section{Lưu trữ media và thanh toán điện tử}
\label{section:3.7}

\subsection{Multer và Cloudinary}

\textbf{Multer} xử lý \texttt{multipart/form-data}; \textbf{Cloudinary} lưu ảnh trên cloud, trả URL công khai~\cite{cloudinary2024}.

\textbf{Yêu cầu liên quan:} ảnh KS/phòng, minh chứng chuyển khoản QR, minh chứng hoàn tiền (UC SJ002, SJ008).

\textbf{Phương án thay thế:} lưu file local (\texttt{/uploads}); AWS S3; Firebase Storage.

\textbf{Lựa chọn:} Cloudinary vì miễn phí tier dev, CDN sẵn, không phụ thuộc ổ đĩa server deploy (Vercel/Railway).

\subsection{VNPay và thanh toán QR}

\textbf{VNPay} là cổng thanh toán phổ biến tại Việt Nam; SDK \texttt{vnpay} tạo URL thanh toán và xác minh chữ ký callback~\cite{vnpay2024}. Song song, \textbf{QR chuyển khoản} cho phép KS tự cấu hình tài khoản; Owner xác nhận thủ công (UC SJ003).

\textbf{Yêu cầu liên quan:} hai phương thức thanh toán trong quy trình đặt phòng (Mục~\ref{subsection:2.2.6}).

\textbf{Phương án thay thế:} MoMo, ZaloPay, PayPal, Stripe.

\textbf{Lựa chọn:} VNPay vì phù hợp thị trường VN và tài liệu sandbox; QR bổ sung cho KS chưa kết nối cổng---linh hoạt vận hành thực tế.

\section{Nền tảng định giá và khuyến mãi}
\label{section:3.8}

\subsection{Tính giá theo đêm và chương trình sale}

Dịch vụ \texttt{salePricingService} áp dụng giảm giá \textbf{theo từng đêm lưu trú}: với mỗi đêm trong khoảng check-in/check-out, chọn sale hiệu lực tốt nhất (phạm vi toàn KS hoặc theo loại phòng), tính \texttt{basePrice}, \texttt{discountAmount}, \texttt{finalAmount} trên server---khách không gửi giá tự tính.

\textbf{Yêu cầu liên quan:} giá xem trước và snapshot khi tạo đơn (UC SJ001); quản lý sale của Owner.

\textbf{Phương án thay thế:} giảm giá cố định cả kỳ; mã coupon nhập tay.

\textbf{Lựa chọn:} tính theo đêm phản ánh đúng nghiệp vụ KS (sale cuối tuần, mùa thấp điểm lẫn lộn trong một kỳ).

\subsection{Gợi ý định giá động (rule-based)}

\texttt{dynamicPricingService} dùng \textbf{luật heuristic}, không phải học máy: hệ số theo \textit{tỷ lệ lấp đầy} phòng, \textit{mùa} (ngày trong năm), \textit{thứ trong tuần} lịch sử; giới hạn biên $\pm$25--35\% so với giá hiện tại~\cite{talluri2004revenue}.

\textbf{Yêu cầu liên quan:} Owner xem gợi ý giá và áp dụng (nếu có trong phạm vi chức năng Chương~\ref{chapter:2}).

\textbf{Phương án thay thế:} hồi quy/ML dự báo nhu cầu; định giá thủ công hoàn toàn.

\textbf{Lựa chọn:} rule-based minh bạch, không cần tập dữ liệu lớn, phù hợp đồ án; ML để dành hướng phát triển.

\section{Công cụ hỗ trợ khác}
\label{section:3.9}

\textbf{ExcelJS} xuất báo cáo Admin. \textbf{dayjs} xử lý ngày giờ (múi giờ VN cho sale). \textbf{React Toastify} hiển thị phản hồi tức thì trên UI.

\section{Tổng kết lựa chọn công nghệ}
\label{section:3.10}

Bảng~\ref{tab:tong-hop-stack} tổng hợp stack chính của StayJourney.

\begin{table}[htbp]
\centering
\caption{Tổng hợp stack công nghệ StayJourney}
\label{tab:tong-hop-stack}
\small
\begin{tabular}{|l|l|}
\hline
\textbf{Tầng} & \textbf{Công nghệ} \\
\hline
Frontend & React 18, Vite, React Router 6, Redux Toolkit, MUI, SCSS, Axios \\
Backend & Node.js, Express 4, Socket.IO \\
Cơ sở dữ liệu & MongoDB, Mongoose \\
Bảo mật & JWT (HttpOnly cookie), bcrypt, 2FA email, RBAC \\
Tích hợp & VNPay, Cloudinary, Nodemailer, node-cron \\
Kiểm tra đầu vào & express-validator \\
\hline
\end{tabular}
\end{table}

Các công nghệ trên tạo nền tảng thống nhất cho thiết kế chi tiết và triển khai ở các chương tiếp theo; mọi lựa chọn đều bám sát yêu cầu nghiệp vụ StayJourney đã phân tích ở Chương~\ref{chapter:2}, ưu tiên tính khả thi triển khai trong phạm vi đồ án tốt nghiệp.

% ---------------------------------------------------------------------------
% Gợi ý: thêm các mục sau vào file tài liệu tham khảo chính (DoAn.tex)
% và dùng \cite{key} như trong chương:
%
% \bibitem{fielding2000rest} R. Fielding, Architectural Styles and the Design
%   of Network-based Software Architectures, PhD thesis, UC Irvine, 2000.
% \bibitem{react2024} Meta, React Documentation, https://react.dev, 2024.
% \bibitem{vite2024} Vite Team, Vite Guide, https://vitejs.dev, 2024.
% \bibitem{redux2024} Redux Team, Redux Toolkit, https://redux-toolkit.js.org, 2024.
% \bibitem{mui2024} MUI, Material UI Documentation, https://mui.com, 2024.
% \bibitem{axios2024} Axios, HTTP Client, https://axios-http.com, 2024.
% \bibitem{nodejs2024} OpenJS Foundation, Node.js Docs, https://nodejs.org, 2024.
% \bibitem{express2024} Express.js, https://expressjs.com, 2024.
% \bibitem{mongodb2024} MongoDB Inc., MongoDB Manual, https://www.mongodb.com/docs, 2024.
% \bibitem{mongoose2024} Mongoose, https://mongoosejs.com, 2024.
% \bibitem{expressvalidator2024} express-validator, https://express-validator.github.io, 2024.
% \bibitem{jones2015jwt} M. Jones et al., RFC 7519: JSON Web Token, IETF, 2015.
% \bibitem{owasp2023session} OWASP, Session Management Cheat Sheet, 2023.
% \bibitem{provos1999bcrypt} N. Provos, D. Mazières, A Future-Adaptable Password Scheme, USENIX, 1999.
% \bibitem{sandhu1996rbac} R. Sandhu et al., Role-Based Access Control Models, IEEE Computer, 1996.
% \bibitem{socketio2024} Socket.IO, Documentation, https://socket.io, 2024.
% \bibitem{nodemailer2024} Nodemailer, https://nodemailer.com, 2024.
% \bibitem{cloudinary2024} Cloudinary, Documentation, https://cloudinary.com/documentation, 2024.
% \bibitem{vnpay2024} VNPay, Tài liệu tích hợp, https://sandbox.vnpayment.vn, 2024.
% \bibitem{talluri2004revenue} K. Talluri, G. van Ryzin, The Theory and Practice of Revenue Management, Springer, 2004.
% ---------------------------------------------------------------------------

\end{document}
